\subsection{Probleme und Anpassungen}
Das erste Problem, das bei der Umsetzung aufgetreten ist, betraf die Frage, über welchen Port der HAProxy erreichbar ist.
Normalerweise wird Port 6443 genutzt, da dieser bei Kubernetes-Nodes der lokale Zugang zu den Kube-APIs ist.
Der Port 6443 ist vom VPN aus jedoch nicht erreichbar, da es sich um einen administrativen Port handelt. 
Das verursacht keine Probleme zwischen VMs im gleichen Netz. Wenn jedoch außerhalb des Netzes über den VPN mit dem Cluster kommuniziert werden soll, ist dieser blockiert.
Deshalb wurde der Port, über den HAProxy erreichbar ist, von 6443 auf 443 geändert.

Ein weiteres Problem ergab sich, da die vorhandenen Worker Nodes im Cluster bei der Reaktivierung ständig selbst auf „notReady” gestellt wurden.
Die Ursache des Problems war, dass die Worker Nodes mit dem Cluster verbunden wurden, aber die Umgebungsvariablen K3S\_TOKEN und K3S\_URL in der Konfiguration fehlten.
Um das Problem zu lösen, wurden die Umgebungsvariablen manuell in der Konfiguration nachgetragen, wodurch die Worker Nodes stabil im Cluster wurden.